% Source: https://artofproblemsolving.com/community/c3h1214053_not_relatively_prime
% Stable identity: aops-topic-1214053-q01
$30030=2\cdot3\cdot5\cdot7\cdot11\cdot13$, so we need every term to be divisible by one of these primes.

\emph{No run of $22$.} Among $22$ consecutive integers, $11$ are odd, and they are $11$ consecutive odd numbers $o_0,o_1,\dots,o_{10}$ (with $o_{i+1}=o_i+2$). Each must be divisible by one of $3,5,7,11,13$. For an odd prime $p$ the multiples of $p$ among them are every $p$-th term, so $3$ divides at most $4$ of them, $5$ at most $3$, $7$ at most $2$, and $11$, $13$ at most $1$ each: at most $4+3+2+1+1=11$. So every bound must be attained and no term may be divisible by two of the primes. But $4$ multiples of $3$ means positions $\{0,3,6,9\}$ or $\{1,4,7,10\}$, and $3$ multiples of $5$ means positions $\{0,5,10\}$; these always share a position, a contradiction.

\emph{A run of $21$.} The integers $9440,9441,\dots,9460$ work: the even ones are divisible by $2$, and the odd ones are
\[9441=3\cdot3147,\ 9443=7\cdot1349,\ 9445=5\cdot1889,\ 9447=3\cdot3149,\ 9449=11\cdot859,\ 9451=13\cdot727,\]
\[9453=3\cdot3151,\ 9455=5\cdot1891,\ 9457=7\cdot1351,\ 9459=3\cdot3153.\]
So the maximum length is $\boxed{21}$.
